%%BLOCK after="Total asphericity tested on final images"
\begin{dcomment}{development}{p97}{A total presheaf whose category of elements is not total}
Grothendieck expects that Proposition~8 has no converse outside the subterminal case, and proposes to look for a totally $W$-aspherical presheaf $H$ whose category of elements is not totally $W$-aspherical. Such a presheaf exists on the smallest category that is not a preorder and not a group. We found no source for the example; it is given without a claim of priority.

\textbf{Proposition (the converse of Proposition~8 fails).}\ours
\textit{Let $U$ be the category with one object $*$ and $\operatorname{End}(*)=\{1,e\}$, $e^2=e$. Let $H$ be the presheaf on $U$ with $H(*)=\{i,h_1,h_2\}$, on which $e$ acts by the constant map with value $i$. Then $H$ is totally $W$-aspherical for every fundamental localizer $W$. The category $U/H$ is not totally $W$-aspherical for $W=W_\infty$, for the homology equivalences $W_R$ with coefficients in a nonzero commutative ring $R$, and for the $n$-equivalences with $n\geq2$.}

\textit{Proof.}
Let $R=U(-,*)$ be the representable presheaf. Its elements are $1$ and $e$, and $e$ acts on them by $x\mapsto xe$, the constant map with value $e$. The endomorphism $\varepsilon:x\mapsto ex$ of $R$ is idempotent with image the subpresheaf $Q=\{e\}$. The presheaf $H$ is the pushout $R\amalg_QR$ of two copies of $R$ along $Q$: the element $e$ is $i$, and the two copies of $1$ are $h_1,h_2$.

Let $p:U\to C$ be any functor to a small category, and put $c=p(*)$, $d=p(e)$, so that $d^2=d$. The functor $p_!$ preserves colimits, sends $R$ to $C(-,c)$, and sends $\varepsilon$ to the endomorphism $\alpha\mapsto d\alpha$. Every functor preserves the splitting of an idempotent. Hence $p_!Q$ is the subpresheaf
\[
 Q_d(v)=\{\alpha:v\to c\mid d\alpha=\alpha\}\subseteq C(v,c),
\]
with its inclusion, and
\[
 p_!H\cong C(-,c)\amalg_{Q_d}C(-,c).
\]
This pushout is formed objectwise along injections. The category of elements $E'=C/p_!H$ is therefore the union of two full sieves $S_1,S_2$, each isomorphic to $C/c$, with $S_1\cap S_2=C/Q_d$. A chain of arrows of $E'$ lies in the sieve that contains its last object, so $NE'=NS_1\cup NS_2$ and $NS_1\cap NS_2=N(C/Q_d)$. Each $S_j$ has the terminal object $(c,\mathrm{id}_c)$. The category $C/Q_d$ has the object $(c,d)$, and the arrows $\alpha:(v,\alpha)\to(c,d)$ are the components of a natural transformation from the identity functor to the constant functor at $(c,d)$: for an arrow $g:(v,\alpha)\to(w,\beta)$ one has $\beta g=\alpha$. Thus $BS_1$, $BS_2$ and their intersection are contractible, and so is $BE'$. Since $p$ was arbitrary, $H$ is totally $W_\infty$-aspherical. Every fundamental localizer contains $W_\infty$ \citep[Th\'eor\`eme~2.2.11]{Cisinski2004Minimal}, which gives the first assertion.

The category $E=U/H$ has the three objects $i,h_1,h_2$ and, besides identities, the arrows
\[
 e_i:i\to i,\qquad u_1:i\to h_1,\qquad u_2:i\to h_2,
 \qquad e_i^2=e_i,\quad u_je_i=u_j.
\]
Let $q:E\to\Psi$ be the functor to the outgoing fork $\Psi=(a\leftarrow z\to b)$ of the commentary at ms.~100 which sends $i,h_1,h_2$ to $z,a,b$ and $e_i$ to the identity. It is bijective on objects and full. For an object $y$ of $\Psi$ with preimage $\tilde y$, every object $(x,\,y\to q(x))$ of $y\backslash E$ receives an arrow from $(\tilde y,\mathrm{id}_y)$, obtained by lifting $y\to q(x)$. So $q$ is final. The inclusion $f:\Psi\to C$ into the four-object category of the same commentary is final and $BC\simeq S^2$. The composite $fq:E\to C$ is a final functor whose target is not aspherical for any of the localizers named. By the final-image criterion above, $U/H$ is not totally aspherical for them.\hfill$\square$

The objects $h_1,h_2$ of $U/H$ have no common successor; this is the property that fails, and it cannot be seen by functors that factor through $U$. For the localizers $W_0$ and $W_1$ (bijections on $\pi_0$; bijections on $\pi_0$ and isomorphisms on all $\pi_1$) the converse of Proposition~8 does hold. A final functor induces a bijection on connected components, and a final image of a $1$-connected category is $1$-connected (the argument with coverings at ms.~101). If $H$ is totally $W_0$-aspherical, the test $p=\mathrm{id}$ shows that $U/H$ is nonempty and connected, and then so is every final image of $U/H$; for $W_1$ the same reasoning applies with $1$-connectedness.
\end{dcomment}

%%BLOCK after="The information lost by taking components"
\begin{dcomment}{development}{p103}{The path category over a contractible base need not be total}
The question that closes this page, whether $\underline{\mathrm{Ch}}(Y;y_0,-)$ is totally $W_\infty$-aspherical when $Y$ is $W_\infty$-aspherical, has a negative answer, and the witness is the category that Grothendieck names in the footnote to Proposition~10.

\textbf{Proposition.}\ours
\textit{Let $Y=\Psi$ be the outgoing fork $a\leftarrow z\to b$ and $y_0=z$. Then $BY$ is contractible, and $P=\underline{\mathrm{Ch}}(Y;z,-)$ is $W_\infty$-aspherical and not totally $W_\infty$-aspherical. It is not totally $W_R$-aspherical for any nonzero commutative ring $R$ either.}

\textit{Proof.}
The object $z$ is initial, so $BY$ is contractible; in particular $Y$ is $1$-connected. Proposition~10 applies: $P$ is $W_\infty$-aspherical, and, as recalled in the preceding commentary, the endpoint functor $f:P\to Y$ satisfies $f_!e_P=e_Y$, that is, every $y\backslash P$ is nonempty and connected. So $f$ is final. The inclusion $g:\Psi\to C$ of the commentary at ms.~100 is final, and $BC\simeq S^2$. Final functors compose, because $(gf)_!e_P=g_!f_!e_P=g_!e_Y=e_C$. The category of elements of $(gf)_!e_P$ is $C$, which is neither $W_\infty$-aspherical nor $W_R$-aspherical.\hfill$\square$

The proof uses only two facts: the endpoint functor of the path category over a $1$-connected base is final, and total asphericity passes to final images. Hence $\underline{\mathrm{Ch}}(Y;y_0,-)$ is totally $W_\infty$-aspherical only if $Y$ is. Since total asphericity of $Y$ is far stronger than asphericity (for a category with finitely many objects it means that $Y$ is filtered; see the commentary on Chapter~II, ms.~36), the path construction does not produce totally aspherical categories from aspherical ones.
\end{dcomment}

%%BLOCK after="The later theory of homotopy cofibrations"
\begin{dcomment}{development}{p111}{The open questions of this page}
Question~a) asks whether a sieve that is a $W$-cofibration, or even a trivial $W$-cofibration, must be $W$-dwy\'erienne, and asserts this when the complement $T$ is the final category. Three small examples separate the notions. In each of them $U=\{u\}$ is a single object, so a functor $p:U\to A$ is the choice of an object $c=p(u)$ of $A$.

\textbf{Proposition (a trivial cofibration that is not dwy\'erienne).}\ours
\textit{Let $X$ have objects $u,t,s$ and nonidentity arrows $\alpha,\beta:u\to t$, $\gamma:t\to s$, $\delta:u\to s$, with $\gamma\alpha=\gamma\beta=\delta$. The sieve $U=\{u\}\hookrightarrow X$ is a trivial $W$-cofibration for every fundamental localizer $W$. It is not $W$-dwy\'erienne when $W$ is $W_\infty$, a homology localizer $W_R$ with $R\neq0$, or the class $W_n$ of $n$-equivalences for some $n\geq0$.}

\textit{Proof.}
The presheaf $v\mapsto\operatorname{Hom}_X(v,t)$ on $U$ has two elements, so $U/t$ is a discrete category with two objects; the test $p=\mathrm{id}$ shows that this presheaf is not totally $W$-aspherical for the localizers named. Let $p:U\to A$ be arbitrary, $c=p(u)$, and $P=X\amalg_UA$. Its objects are those of $A$ and $t,s$; for $v\in A$,
\[
 P(v,t)=A(v,c)\times\{\alpha,\beta\},\qquad P(v,s)=A(v,c),\qquad P(t,s)=\{\gamma\},
\]
composition with $\gamma$ forgets the second coordinate, and there are no arrows from $t$ or $s$ to $A$. Let $j:A\to P$ be the inclusion and $r:P\to A$ the functor that is the identity on $A$, sends $t,s$ to $c$, projects the two crossing sets to $A(v,c)$ and sends $\gamma$ to $\mathrm{id}_c$. Let $F:P\to P$ be the identity on $A$, send $t$ and $s$ to $s$, each crossing arrow to its composite into $s$, and $\gamma$ to $\mathrm{id}_s$. There are natural transformations
\[
 \mathrm{id}_P\Longrightarrow F\Longleftarrow jr,
\]
both the identity on $A$; the first has components $\gamma$ at $t$ and $\mathrm{id}_s$ at $s$, the second has at $t$ and at $s$ the arrow $c\to s$ corresponding to $\mathrm{id}_c$. Naturality at a crossing arrow $(g,\alpha)$ or $(g,\beta)$ is the relation $\gamma\alpha=\gamma\beta=\delta$. Since $rj=\mathrm{id}_A$, the map $Bj$ is a homotopy equivalence. So every cobase change of $U\to X$ is in $W_\infty\subseteq W$.\hfill$\square$

This answers the first sentence of question~a) negatively; the footnote's stronger hope fails with it, since a Dwyer immersion is dwy\'erienne.

\textbf{Proposition (the case $T\simeq e$).}\ours
\textit{Let $X_0$ have objects $u,t$ and exactly two nonidentity arrows $\alpha,\beta:u\to t$, and let $U=\{u\}$, $T=\{t\}$. For $W=W_\infty$ and for $W=W_R$, the sieve $U\hookrightarrow X_0$ and every sieve $U\hookrightarrow X_1$ deduced from it by a change of complement $T_1\to T$ (Proposition~7~b)) is a $W$-cofibration, and none of them with $T_1\neq\varnothing$ is $W$-dwy\'erienne.}

\textit{Proof.}
The category $X_1$ has the objects $u$ and those of $T_1$, the arrows of $T_1$, and two arrows $\alpha_x,\beta_x:u\to x$ for every $x\in T_1$, compatible with the arrows of $T_1$. Each crossing presheaf has two elements, so no $U/x$ is connected. For $p:U\to A$ with $c=p(u)$, the pushout $P_A=X_1\amalg_UA$ has $P_A(v,x)=A(v,c)\times\{\alpha,\beta\}$ for $v\in A$, $x\in T_1$. A chain of arrows of $P_A$ either lies in $A$, or lies in $T_1$, or consists of a chain in $A$, a crossing arrow $(g,\lambda)$ with $\lambda\in\{\alpha,\beta\}$, and a chain in $T_1$; the chain in $A$ together with $g$ is a chain of $A/c$. Write $J=N(A/c)*NT_1$ for the join, and $J_\alpha\cup J_\beta$ for two copies of $J$ glued along $NT_1$. Then $NP_A$ is the pushout of simplicial sets
\[
 NA\longleftarrow N(A/c)\amalg N(A/c)\hookrightarrow J_\alpha\cup J_\beta ,
\]
in which the left map forgets the arrow to $c$ and the right map is a monomorphism; it is therefore a homotopy pushout. A functor $h:A\to A'$ under $U$ induces a map of these diagrams. The categories $A/c$ and $A'/h(c)$ have terminal objects, so their nerves and the joins $J$ are contractible, and the maps on $N(A/c)\amalg N(A/c)$ and on $J_\alpha\cup J_\beta$ are weak equivalences whatever $h$ is. If $h\in W_\infty$, the map $NP_A\to NP_{A'}$ is then a weak equivalence. For $W_R$ one uses the Mayer--Vietoris sequences of the same pushouts.\hfill$\square$

So the assertion ``je sais que c'est le cas si on suppose $T$ isomorphe \`a la cat\'egorie finale $e$'' is false for $W$-cofibrations, and with it the equivalence stated at the end of question~a), also under the axiom W(7), which the localizers $W_R$ satisfy. The corollary of Proposition~7 concerns membership in $W$ of one inclusion, not the cofibration property. What it does give is the statement for \emph{trivial} cofibrations and homology:

\textbf{Proposition.}\ours
\textit{Let $U\hookrightarrow X$ be a sieve whose complement is one object $t$ with only its identity, and suppose that $H=\operatorname{Hom}_X(-,t)|_U$ is not empty. For a nonzero commutative ring $R$, the inclusion is a trivial $W_R$-cofibration if and only if $H$ is totally $W_R$-aspherical, that is, if and only if it is $W_R$-dwy\'erienne.}

\textit{Proof.}
After a cobase change $p:U\to A$ the complement is still $\{t\}$ and the crossing presheaf is $p_!H$. A chain of arrows of $X\amalg_UA$ that does not lie in $A$ is a chain of $A/p_!H$ followed by $t$, so $N(X\amalg_UA)$ is $NA$ with the cone on $N(A/p_!H)$ attached along the forgetful map, and $B(X\amalg_UA)$ is the mapping cone of $B(A/p_!H)\to BA$. The inclusion of $BA$ is an $R$-homology isomorphism exactly when the reduced $R$-homology of $A/p_!H$ vanishes, that is, when $A/p_!H$ has the $R$-homology of a point. Quantifying over $p$ gives the statement.\hfill$\square$

For $W_\infty$ the same computation shows that a one-point sieve is a trivial $W_\infty$-cofibration if and only if every $A/p_!H$ is nonempty and integrally acyclic and the cone attachment does not change fundamental groups. Whether this forces $H$ to be totally $W_\infty$-aspherical is not decided here: a counterexample would be a presheaf $H$ all of whose direct images have acyclic categories of elements, one of them with a nontrivial (necessarily perfect) fundamental group.
\end{dcomment}

%%BLOCK after="Dwyer attachments and higher categorical pushouts"
\begin{dcomment}{development}{p114}{Composition of aspherical functors and of dwy\'erienne immersions}
Grothendieck doubts that strictly or totally $W$-aspherical functors, and the two classes of dwy\'erienne immersions, are stable under composition. For functors the doubt is unfounded; for immersions it is justified for some localizers and not for others. The statements below were not found in the sources consulted; no priority is claimed.

\textbf{Theorem (functors compose).}\ours
\textit{Let $W$ be a fundamental localizer.}
\begin{enumerate}[label=(\alph*)]
\item \textit{If $f:X\to Y$ is strictly $W$-aspherical and $H$ is a totally $W$-aspherical presheaf on $Y$, then $f^*H$ is a totally $W$-aspherical presheaf on $X$. Consequently strictly $W$-aspherical functors are stable under composition.}
\item \textit{Totally $W$-aspherical functors are stable under composition.}
\item \textit{If $B$ is a totally $W$-aspherical category and $D:B\to\mathrm{Cat}$ a functor with totally $W$-aspherical values, the cofibred category $\int_BD$ is totally $W$-aspherical.}
\end{enumerate}

\textit{Proof.}
(a) Let $p:X\to X'$ be arbitrary. For $y\in Y$ put $Q_y=p_!\operatorname{Hom}_Y(f(-),y)$, a presheaf on $X'$, covariant in $y$. Each $Q_y$ is totally $W$-aspherical, by the hypothesis on $f$ and the lemma of ms.~112. Let $Z$ be the category with full subcategories $X'$ and $Y$, with $\operatorname{Hom}_Z(x',y)=Q_y(x')$ and no arrows from $Y$ to $X'$; let $j:Y\hookrightarrow Z$ be the inclusion and $F=j_!H$. Then $F|_Y=H$ and
\[
 F|_{X'}=\int^{y\in Y}Q_y\times H(y)\cong p_!\Bigl(\int^{y}\operatorname{Hom}_Y(f(-),y)\times H(y)\Bigr)=p_!f^*H=:K .
\]
Since $H$ is total, $Z/F$ is $W$-aspherical. The inclusion $l:X'/K\hookrightarrow Z/F$ is a sieve. Its comma category over an object $(x',k)$ of $X'/K$ has a terminal object. Its comma category over $(y,h)$, $h\in H(y)$, has as objects the pairs $(x',q)$ with $q\in Q_y(x')$ (the element of $K(x')$ is then the class of $(q,h)$) and as arrows the arrows of $X'$ that respect $q$; it is $X'/Q_y$, which is $W$-aspherical. By the relative form of Theorem~A that defines a fundamental localizer, $l\in W$, and so $X'/K$ is $W$-aspherical. As $p$ was arbitrary, $f^*H$ is total. If $g:Y\to Z_0$ is also strictly aspherical, apply this to $H=\operatorname{Hom}_{Z_0}(g(-),z)$, whose restriction along $f$ is $\operatorname{Hom}_{Z_0}(gf(-),z)$.

(b) Let $f:X\to Y$ and $g:Y\to Z_0$ be totally aspherical and $z\in Z_0$. The functor $f_z:X/z\to Y/z$, $(x,\gamma)\mapsto(f(x),\gamma)$, has comma category over $(y,\beta)$ canonically isomorphic to $X/y$, because the arrow to $z$ is determined as $\beta g(h)$ by $h:f(x)\to y$. So $f_z$ is totally aspherical, hence strictly aspherical by Proposition~8. The final presheaf of $Y/z$ is total because $g$ is totally aspherical; by (a) its restriction, the final presheaf of $X/z$, is total. That is, $X/z$ is a totally $W$-aspherical category.

(c) Put $E=\int_BD$, with objects $(b,x)$, $x\in D(b)$, and projection $\pi:E\to B$. For $b\in B$ the functor $D(b)\to E/b$, $x\mapsto(b,x,\mathrm{id}_b)$, has the left adjoint $(b',x,\beta)\mapsto D(\beta)(x)$. A right adjoint is final. So $E/b$ is a final image of $D(b)$ and is total; $\pi$ is a totally aspherical functor, as is $B\to e$, and (b) applies.\hfill$\square$

\textbf{Proposition (immersions: two counterexamples).}\ours
\textit{For $W=W_0$ and for $W=W_1$ neither the $W$-dwy\'erienne nor the totally $W$-dwy\'erienne sieves are stable under composition.}

\textit{Proof.}
Let $V$ be the discrete category on two objects for $W_0$, and the category with two objects and two parallel arrows for $W_1$; $BV$ is $S^0$, respectively $S^1$. Let $B=V^{\triangleleft}$ be $V$ with an initial object $o$ adjoined. Let $X$ consist of $B$ and of a second copy $V_0$ of $V$, with
\[
 X(v_0,b)=B(v,b),\qquad X(b,v_0)=\varnothing\qquad(v\in V,\ b\in B),
\]
and composition induced by $B$. The sieve $V_0\hookrightarrow X$ generates the cosieve $V_0\cup V$, on which the inclusion of $V_0$ has the right adjoint $v\mapsto v_0$; it is a Dwyer immersion. The functor $X\to B$ forgetting the copy is left adjoint to the inclusion $B\subset X$, so $BX\simeq BB$ is contractible. A final functor is bijective on $\pi_0$ and surjective on $\pi_1$; hence $X$ is totally $W_0$- and totally $W_1$-aspherical. Let $Y=X^{\triangleright}$ be $X$ with a terminal object $t$ adjoined. The sieve $X\hookrightarrow Y$ is totally $W$-dwy\'erienne, its only boundary category being $X/t=X$. For the composite sieve $V_0\hookrightarrow Y$ the boundary category at $t$ is $V_0/t\simeq V$ and the boundary presheaf is the final presheaf of $V_0$; neither is $W$-aspherical.\hfill$\square$

The targets $Y$ have six objects and $15$, respectively $21$, arrows.

\textbf{Theorem (total immersions for $W_\infty$).}\ours
\textit{For $W=W_\infty$ the totally $W$-dwy\'erienne sieves are stable under composition.}

\textit{Proof.}
We use the criterion of the commentary on Chapter~II, ms.~36: a category $A$ is totally $W_\infty$-aspherical if and only if it is nonempty and connected and every connected component of the category of elements of every functor $A\to\mathrm{Set}$ is weakly contractible; and then those components are themselves totally $W_\infty$-aspherical and any two objects of $A$ have a common successor. A nonempty cosieve $S$ of such an $A$ is the category of elements of the functor with value a point on $S$ and empty elsewhere; it is connected, because two of its objects have a common successor in $A$, which lies in $S$. So $S$ is totally $W_\infty$-aspherical.

Let $U\subset X\subset Y$ be totally $W_\infty$-dwy\'erienne sieves and $y\in Y$ with $U/y\neq\varnothing$. Then $X/y$ is total (it has a terminal object if $y\in X$). The objects $(x,\beta)$ of $X/y$ with $U/x\neq\varnothing$ form a nonempty cosieve $S_y$, total by the preceding paragraph, and the functor $(x,\beta)\mapsto U/x$ on $S_y$ has total values. By part~(c) of the theorem its cofibred category $G$ is total. The functor $G\to U/y$ that composes $u\to x\to y$ has the left adjoint $(u,\gamma)\mapsto((u,\gamma),(u,\mathrm{id}_u))$, so it is final, and $U/y$ is total.\hfill$\square$

The same proof works for the homology localizers $W_R$, with the homological form of the criterion ($R$-acyclic in place of weakly contractible). For the strict class, the $W_R$-dwy\'erienne sieves are also stable under composition: the restriction of a totally $W_R$-aspherical presheaf to a cosieve on which it is nonempty is totally $W_R$-aspherical, after which part~(a) applies; the proof of the restriction statement is not reproduced here. For $W_\infty$-dwy\'erienne sieves the question is undecided. It comes down to the question left open in the commentary at ms.~111: does a presheaf all of whose direct images have integrally acyclic categories of elements have only weakly contractible ones?
\end{dcomment}

\begin{dcomment}{development}{p114}{The endpoint product is not a cofibration}
The class asked for in b) does not exist for $W_\infty$: the inclusion $\alpha$ built from the simplest sieve is not a $W_\infty$-cofibration at all. This is stronger than the observation that $\alpha$ is not a Dwyer immersion.

\textbf{Theorem.}\ours
\textit{Let $n\geq2$, $I=\Phi_n$, $\varepsilon=\{0,2n\}\subset I$, and let $\alpha:A\hookrightarrow P$ be the inclusion of $A=(\{0\}\times I)\cup([1]\times\varepsilon)$ in $P=[1]\times I$, deduced from $\{0\}\hookrightarrow[1]$. Let $C$ be the four-object category of the commentary at ms.~100, of the homotopy type of $S^2$.}
\begin{enumerate}[label=(\alph*)]
\item \textit{If $W$ is a fundamental localizer for which a category of the homotopy type of $S^3$ is not $W$-aspherical (for instance $W_\infty$, $W_R$ with $R\neq0$, or the $m$-equivalences with $m\geq3$), then $\alpha$ is not a $W$-cofibration.}
\item \textit{If $C$ is not $W$-aspherical (which includes the $2$-equivalences), then $\alpha$, which belongs to $W$, is not a trivial $W$-cofibration.}
\end{enumerate}

\textit{Proof.}
Write $l_k=(0,k)$ and $u_k=(1,k)$. The objects of $A$ are $l_0,\dots,l_{2n},u_0,u_{2n}$; its nonidentity arrows are those of the lower copy of $I$ and $l_0\to u_0$, $l_{2n}\to u_{2n}$, and no two of them are composable. Both $A$ and $P$ have contractible classifying spaces, so $\alpha\in W_\infty\subseteq W$.

Recall that $C$ has objects $x,z,a,b$ and arrows $\alpha',\beta':x\to z$, $r:z\to a$, $s:z\to b$ with $r\alpha'=r\beta'$, $s\alpha'=s\beta'$. Let $D$ be either $C$ or the category $C^{\triangleright}$ obtained by adjoining a terminal object $\omega$. Define $p:A\to D$ by
\[
 p(l_0)=z,\qquad p(u_0)=b,\qquad p(l_k)=a\ (1\leq k\leq2n),\qquad p(u_{2n})=a,
\]
sending $l_0\to l_1$ to $r$, $l_0\to u_0$ to $s$ and all other arrows to $\mathrm{id}_a$. Let $Q=P\amalg_AD$. Since $A$ is a sieve in $P$, the category $Q$ consists of $D$, of the objects $u_1,\dots,u_{2n-1}$ with the arrows of the upper copy of $I$ between them, and of the crossing sets $Q(v,u_k)=(p_!H_k)(v)$, where $H_k=\operatorname{Hom}_P(-,u_k)|_A$ (Theorem~1).

The presheaf $H_1$ is the final presheaf of the sieve $A_1=\{l_0,l_1,l_2,u_0\}$ of $A$ and is empty elsewhere. So $(p_!H_1)(v)=\pi_0(v\backslash A_1)$, the comma category being formed with the restriction $A_1\to D$ of $p$. For $v=z$ it is $A_1$ itself, a tree; for $v=a$ it is $l_2\to l_1$; for $v=b$ it is the point $u_0$; for $v=x$ it has the five objects $(l_0,\alpha'),(l_0,\beta'),(l_1,r\alpha'),(l_2,r\alpha'),(u_0,s\alpha')$ and is connected; for $v=\omega$ it is empty. Hence the object $u_1$ receives exactly one arrow from each object of $C$ and none from $\omega$. For $2\leq k\leq2n-1$ the category $A/u_k$ is nonempty and connected and $p$ is constant on it with value $a$, so $p_!H_k=D(-,a)$; the arrows between $u_2,\dots,u_{2n-1}$ induce identities of $D(-,a)$, and $u_2\to u_1$ induces the unique map $D(-,a)\to p_!H_1$.

Let $R\subset Q$ be the full subcategory on $D$ and $u_1$. The functor $\rho:Q\to R$ that is the identity on $R$, sends $u_k$ to $a$ for $k\geq2$, is the identity on the crossing sets $D(v,a)$, sends the arrows between these $u_k$ to $\mathrm{id}_a$ and $u_2\to u_1$ to the arrow $a\to u_1$, satisfies $\rho|_R=\mathrm{id}_R$; and the arrows $a\to u_k$ corresponding to $\mathrm{id}_a$ are the components of a natural transformation from $\rho$, followed by the inclusion of $R$, to $\mathrm{id}_Q$. So $BQ\simeq BR$.

(a) Take $D=C^{\triangleright}$. Then $p$ is a functor between categories with contractible classifying spaces, so $p\in W_\infty\subseteq W$; it is a morphism of $A\backslash\mathrm{Cat}$ from $\mathrm{id}_A$ to $p$, and its image under cobase change along $\alpha$ is $P\to Q$. The category $R$ is $C$ with the two terminal objects $\omega$ and $u_1$ adjoined, with no arrow between them: its nerve is the union of two cones on $NC$, so $BR\simeq\Sigma BC\simeq S^3$. Thus $P$ is $W$-aspherical and $Q$ is not, and $P\to Q$ is not in $W$.

(b) Take $D=C$. Then $R$ is $C$ with one terminal object adjoined and $BQ\simeq BR$ is contractible. The cobase change $C\to Q$ of $\alpha$ along $p$ goes from a category that is not $W$-aspherical to one that is, so it is not in $W$.\hfill$\square$

For $n=2$ the category $Q$ of (a) has $8$ objects and $30$ arrows. Part~(b) also excludes the cofibration property whenever a $W$-cofibration that belongs to $W$ is a trivial $W$-cofibration. For $W_\infty$ this is \citet[Lemma~1.10]{BataninBerger2017}, applied to Thomason's model structure, which is left proper; there it adds nothing to (a).

\textbf{Corollary (no class of the kind requested).}\ours
\textit{Let $W$ be as in (a) and $n\geq2$. If a class $\mathcal K$ of sieves consists of $W$-cofibrations and contains, with $i:U\hookrightarrow X$, the inclusion of $(U\times I)\cup(X\times\varepsilon)$ in $X\times I$, then every member of $\mathcal K$ is an isomorphism. The same holds for $W$ as in (b) and classes of trivial $W$-cofibrations.}

\textit{Proof.}
$W$-cofibrations, and trivial $W$-cofibrations, are stable under retracts in the category of arrows, because a retract of a functor in $W$ is in $W$. Let $i:U\hookrightarrow X$ be in $\mathcal K$. If there is an arrow $u\to x$ with $u\in U$, $x\notin U$, then $\{0\}\hookrightarrow[1]$ is a retract of $i$ (by the functor $X\to[1]$ with fibre $U$ over $0$, and the functor $[1]\to X$ choosing that arrow), the endpoint product $\alpha$ of $\{0\}\hookrightarrow[1]$ is a retract of that of $i$, and the theorem is contradicted. If there is no such arrow, then $X=U\amalg V$; if $V\neq\varnothing$, the endpoint product of $i$ is the inclusion of $(U\times I)\amalg(V\times\varepsilon)$ in $(U\amalg V)\times I$, which is in $\mathcal K$ and has the crossing arrow $(v,0)\to(v,1)$, so the first case applies to it.\hfill$\square$

At the other extreme, for $W_0$ every functor is a $W_0$-cofibration, because $\pi_0:\mathrm{Cat}\to\mathrm{Set}$ preserves pushouts, and all sieves form a class with the two properties. The same is true for $W_1$ and for every fundamental localizer containing $W_1$: the pushout along a functor injective on objects has the same fundamental groupoid as the cofibred category of the span; this proof is not reproduced here. By \citet[Proposition~9.2.17]{Cisinski2006} a fundamental localizer contains $W_1$ exactly when a category of the homotopy type of $S^2$ is aspherical for it. So for the localizers of homotopy theory proper the endpoint product has to be replaced by a homotopy-invariant construction, which is what the integration of the following pages provides.
\end{dcomment}
